% =====================================================================
%  Module      : Cloud of Things
%  Projet      : ColdTrace — Supervision intelligente de la chaîne du froid
%  Document    : Cahier des charges (Scope Statement)
%  Équipe      : ColdTrace Team
%  Création    : 23/10/2026   Dernière modification : 24/10/2026
%  Version     : 1.0
% =====================================================================
\documentclass[11pt,a4paper]{article}

% ---------- Encodage, langue, polices ----------
\usepackage[utf8]{inputenc}
\usepackage[T1]{fontenc}
\usepackage[french]{babel}
\usepackage{lmodern}
\usepackage{microtype}

% ---------- Mise en page ----------
\usepackage[a4paper,margin=2.5cm]{geometry}
\usepackage{titlesec}
\usepackage{graphicx}
\usepackage{tabularx}
\usepackage{booktabs}
\usepackage{enumitem}
\usepackage{needspace}
\usepackage{xcolor}
\usepackage{tikz}
\usetikzlibrary{positioning,arrows.meta,calc}
\usepackage{pgfgantt}
\usepackage[hidelinks]{hyperref}

\setlength{\parindent}{0pt}
\setlength{\parskip}{0.5em}
\setlist{itemsep=1pt,topsep=3pt}
\newcolumntype{L}{>{\raggedright\arraybackslash}X}

% ---------- Couleurs ----------
\definecolor{ctblue}{HTML}{1F3A5F}
\definecolor{ctcyan}{HTML}{2A9DB8}
\definecolor{ctlight}{HTML}{EAF4F8}
\definecolor{ctgrey}{HTML}{5C6670}

% ---------- Titres : seuls les grands titres sont en couleur ----------
\let\oldsection\section
\renewcommand{\section}{\Needspace{10\baselineskip}\oldsection}
\titleformat{\section}{\Large\bfseries\color{ctblue}}{\thesection}{0.7em}{}
\titleformat{\subsection}{\large\bfseries}{\thesubsection}{0.7em}{}

% ---------- Tableaux (sans couleur) ----------
\newcommand{\tete}[1]{\textbf{#1}}

% ---------- Encadré (texte simple) ----------
\newcommand{\encadre}[1]{\par\textit{#1}\par}

% ---------- Boîtes du diagramme de déploiement ----------
\newcommand{\artbox}[1]{\fcolorbox{ctgrey}{white}{\parbox{2.7cm}{\centering\scriptsize #1}}}
\newcommand{\envbox}[1]{\fcolorbox{ctcyan}{ctlight}{\parbox{3.0cm}{\centering\scriptsize #1}}}

\begin{document}

% =====================================================================
% PAGE DE GARDE
% =====================================================================
\begin{titlepage}
\centering
\vspace*{1.5cm}

% ---- Logo SUP'COM : déposer le fichier logo_supcom.png à côté du .tex ----
\IfFileExists{logo_supcom.png}{%
  \includegraphics[width=4cm]{logo_supcom.png}%
}{%
  \framebox[4cm][c]{\rule{0pt}{1.5cm}\raisebox{0.6cm}{\small Logo SUP'COM}}%
}

\vspace{1cm}
{\Large\bfseries École Supérieure des Communications de Tunis}

\vspace{1.5cm}
\rule{\textwidth}{0.6pt}\\[0.35cm]
{\huge Cahier des charges}\\[0.2cm]
\rule{\textwidth}{0.6pt}

\vspace{1.2cm}
{\large ColdTrace : Supervision de la chaîne du froid}\\[0.2cm]
Projet Cloud of Things

\vspace{1.2cm}
\textit{Réalisé par :}\\
Khalil Hedhili, Oussema Kard, Yessine Abdedayem
% , Prénom Nom   % <-- ajouter un membre ici

\vspace{0.8cm}
\textit{Encadré par :}\\
Dr. Ing. Mohamed-Bécha Kaâniche

\vfill
\textit{Année universitaire :}\\
\textbf{2026-2027}

\vspace{0.5cm}
{\small Dernière modification : 24 octobre 2026 \quad|\quad Last modification: October 24, 2026}
\end{titlepage}

\tableofcontents
\newpage

% =====================================================================
\section{Concept}
% =====================================================================

\subsection{Contexte}
Les produits laitiers, les produits surgelés et de nombreux médicaments doivent rester
dans une plage de température contrôlée, de l'entrepôt jusqu'au client. Beaucoup d'acteurs
de la chaîne du froid utilisent encore des relevés manuels ou des enregistreurs consultés
après coup. Une rupture de température peut donc n'être découverte qu'à la livraison,
lorsque la marchandise est déjà perdue.

\subsection{Problématique}
\encadre{Comment détecter rapidement une anomalie de température ou d'humidité, alerter
immédiatement la bonne personne et garder une trace exploitable de l'incident ?}

\subsection{Ambitions}
Développer une solution \textbf{Cloud of Things de supervision de la chaîne du froid}
utilisant \textbf{Jakarta EE 11 \& Node-RED / MQTT}. Des capteurs connectés mesurent la
température, l'humidité et l'ouverture des portes ; une passerelle IoT transmet les mesures
au cloud, qui les stocke, les analyse et déclenche les alertes.

La solution fournit :
\begin{itemize}
  \item la \textbf{visualisation des données en temps réel} ;
  \item l'\textbf{historique}, des \textbf{analytiques} et des \textbf{indicateurs de performance} :
        temps passé hors plage, nombre d'incidents, délai moyen de réaction à une alerte ;
  \item des \textbf{alertes et notifications} pour chaque anomalie ;
  \item une \textbf{interface de surveillance et de configuration}.
\end{itemize}

Deux applications \textbf{PWA} sont développées :
\begin{itemize}
  \item \textbf{front-office} : tableau de bord temps réel, historique, réception et acquittement des alertes ;
  \item \textbf{back-office} : gestion des capteurs, des sites, des seuils et des utilisateurs.
\end{itemize}
La prédiction n'est pas retenue en version 1 : les alertes reposent sur des règles
mesurables. Elle pourra être étudiée si le temps et les données le permettent.

% =====================================================================
\section{Clients, besoins et objectifs}
% =====================================================================

\subsection{Clients et utilisateurs}
\begin{tabularx}{\textwidth}{@{}l L@{}}
\toprule
\tete{Client cible} & \tete{Usage principal} \\
\midrule
Distributeurs alimentaires, laiteries & Suivi des chambres froides et des quais \\
Pharmacies, grossistes & Conservation des médicaments thermosensibles \\
Transporteurs frigorifiques & Surveillance des camions pendant la livraison \\
\bottomrule
\end{tabularx}

Trois profils utilisent la solution : le \textbf{responsable qualité} (alertes, analyse,
rapports), le \textbf{chauffeur ou opérateur} (alertes de son véhicule ou de sa zone) et
l'\textbf{administrateur} (configuration dans le back-office).

\subsection{Besoins fonctionnels}
\begin{tabularx}{\textwidth}{@{}l L l@{}}
\toprule
\tete{ID} & \tete{Besoin} & \tete{Priorité} \\
\midrule
EF-01 & Mesurer en continu la température, l'humidité et l'ouverture de porte & Haute \\
EF-02 & Afficher les mesures en temps réel & Haute \\
EF-03 & Alerter en cas de dépassement de seuil : froid 0--4\,°C, surgelé $\le -18$\,°C, humidité 50--70\,\% & Haute \\
EF-04 & Notifier le responsable qualité et le chauffeur concerné & Haute \\
EF-05 & Acquitter et commenter une alerte & Moyenne \\
EF-06 & Conserver l'historique des mesures et des incidents & Haute \\
EF-07 & Afficher les KPI et exporter un rapport de conformité & Moyenne \\
EF-08 & Gérer capteurs, sites, seuils et utilisateurs & Haute \\
\bottomrule
\end{tabularx}

\subsection{Besoins non fonctionnels}
\begin{itemize}
  \item \textbf{Réactivité} : une alerte s'affiche moins de 10 secondes après la mesure ;
  \item \textbf{Sécurité} : MQTTS et HTTPS (TLS), accès par rôle, secrets hors des dépôts ;
  \item \textbf{Fiabilité} : la passerelle garde les mesures en mémoire en cas de coupure réseau ;
  \item \textbf{Ergonomie} : PWA installables sur mobile et simples à utiliser.
\end{itemize}

\subsection{Occasions d'utilisation}
\begin{itemize}
  \item \textbf{Transport} : surveillance pendant la livraison, alerte au chauffeur ;
  \item \textbf{Stockage} : suivi permanent des chambres froides ;
  \item \textbf{Audit qualité} : consultation de l'historique et des incidents lors d'un contrôle.
\end{itemize}

\subsection{Limites (hors périmètre)}
La solution \textbf{surveille et informe} ; elle ne remplace ni l'équipement frigorifique
ni les procédures de conformité. Sont exclus : le réglage du groupe froid, la gestion
complète de flotte, la facturation, la certification réglementaire, les protocoles autres
que MQTT (CoAP, AMQP) et la prédiction par IA.

% =====================================================================
\section{Technologies choisies}\label{sec:techno}
% =====================================================================
Les technologies retenues couvrent la partie logicielle et la partie matérielle du projet.

\subsection{Partie logicielle}

\textbf{Middleware}
\begin{itemize}
  \item \textbf{WildFly} : serveur d'applications open source compatible Jakarta EE 11, qui héberge l'API REST.
  \item \textbf{Mosquitto} : broker MQTT léger qui reçoit les mesures de la passerelle et les distribue au serveur ; configuré en MQTTS (TLS, port 8883).
  \item \textbf{Jakarta REST (JAX-RS)} : API Java pour exposer les services web REST consommés par les PWA.
\end{itemize}

\textbf{Back-end}
\begin{itemize}
  \item \textbf{Jakarta EE 11} : CDI, Jakarta REST, Jakarta WebSocket (envoi des mesures et des alertes en temps réel) et Jakarta Security (authentification par rôle).
  \item \textbf{MongoDB} : base de données NoSQL orientée documents, adaptée aux séries de mesures, aux incidents et à la configuration.
  \item \textbf{MQTT} : protocole publish/subscribe léger utilisé entre la passerelle et le cloud.
\end{itemize}

\textbf{Front-end}
\begin{itemize}
  \item \textbf{PWA (HTML5, CSS3, JavaScript)} : deux applications web installables sur mobile, avec mode hors ligne (Service Worker) et notifications push.
  \item \textbf{Chart.js} : graphiques temps réel et historiques du tableau de bord.
\end{itemize}

\textbf{Intégration IoT}
\begin{itemize}
  \item \textbf{Node-RED} : outil de programmation par flux installé sur la passerelle ; il collecte les mesures des capteurs, les met en forme, les garde en mémoire en cas de coupure et les publie sur le broker.
  \item \textbf{Arduino IDE / C++} : programmation du firmware des ESP32.
\end{itemize}

\textbf{Outils}
\begin{itemize}
  \item \textbf{GitHub} : versionnement du code et déploiement sur le cloud (\texttt{git clone} puis \texttt{git pull}).
  \item \textbf{Postman} : collection de tests de l'API REST.
  \item \textbf{Cloud IaaS (Ubuntu Linux)} : hébergement du broker, du serveur et de la base de données.
\end{itemize}

\subsection{Partie matérielle}
\begin{tabularx}{\textwidth}{@{}l L@{}}
\toprule
\tete{Composant} & \tete{Rôle} \\
\midrule
Raspberry Pi 4 & Passerelle IoT (Node-RED) \\
ESP32 (1 à 2) & Microcontrôleurs Wi-Fi reliés aux capteurs \\
DS18B20 (sonde étanche) & Température, y compris en surgelé (jusqu'à $-55$\,°C) \\
DHT22 & Température et humidité de l'air \\
Capteur magnétique de porte (reed) & Détection d'ouverture de porte \\
Buzzer et LED & Actionneurs : alerte locale en cas d'anomalie \\
Capteurs simulés & Script qui complète la flotte réelle pour les tests \\
\bottomrule
\end{tabularx}

% =====================================================================
\section{Modèle commercial}
% =====================================================================

\subsection{Stratégie de commercialisation}
\begin{itemize}
  \item \textbf{Étude de marché} : les PME de la chaîne du froid (laiteries, distributeurs,
        pharmacies, transporteurs) sont peu équipées en supervision connectée ; les solutions
        existantes sont souvent consultées après coup ou trop coûteuses pour elles.
  \item \textbf{Différenciation} : traçabilité complète de chaque incident (horodatage,
        acquittement, commentaire), utile lors d'un audit, et abonnement sans investissement logiciel.
  \item \textbf{Jalons} : bêta-test avec un client pilote au T1 2027, lancement au T2 2027.
\end{itemize}

\subsection{Politique marketing (4 P)}
\begin{tabularx}{\textwidth}{@{}l L@{}}
\toprule
\tete{Politique} & \tete{Choix} \\
\midrule
Produit & Kit capteurs ESP32, passerelle Raspberry Pi, plateforme cloud et deux PWA \\
Prix & Abonnement de 8 à 15\,€ par capteur et par mois selon le volume ; matériel vendu ou loué \\
Communication & Démonstration, vidéo, dossier de presse, LinkedIn et salons professionnels \\
Distribution & Vente directe aux PME et partenariats avec des frigoristes et logisticiens \\
\bottomrule
\end{tabularx}

\textit{Le prix, le client pilote et les dates de lancement sont des hypothèses de travail à valider.}

% =====================================================================
\section{Livrables}
% =====================================================================
\begin{tabularx}{\textwidth}{@{}l L@{}}
\toprule
\tete{Livrable} & \tete{Contenu} \\
\midrule
Cahier des charges & Le présent document \\
Cahier conceptuel & Besoins, cas d'utilisation, diagrammes UML, patrons de conception \\
Codes sources & Dépôts GitHub : API REST, PWA front-office et back-office, passerelle IoT \\
Déploiement & API REST sur Cloud IaaS, testable avec Postman ; PWA installées sur nos téléphones \\
Documentation & Documentation technique et manuel utilisateur \\
Communication & Supports marketing, dossier de presse, campagne publicitaire \\
Vidéos & Une vidéo de présentation de moins de 5 minutes par livrable \\
\bottomrule
\end{tabularx}

% =====================================================================
\section{Contraintes}
% =====================================================================

\subsection{Temps, argent, objectifs}
\begin{itemize}
  \item \textbf{Temps} : livraison finale fin janvier 2027, en parallèle des autres modules ;
  \item \textbf{Argent} : matériel limité (1 Raspberry Pi, 1 à 2 ESP32), complété par des capteurs simulés ;
  \item \textbf{Objectifs} : priorité aux alertes fiables et au temps réel.
\end{itemize}

\subsection{Méthodologie : eXtreme Programming}
Itérations de deux semaines terminées par une démonstration, \textit{user stories} issues
des besoins EF-01 à EF-08, programmation en binôme sur les parties critiques, tests continus
et intégration régulière sur GitHub. Le déploiement sur le cloud se fait uniquement depuis
GitHub (\texttt{git clone} puis \texttt{git pull}).

\subsection{Contraintes technologiques}
Le module impose Jakarta EE 11 pour le serveur, Node-RED et MQTT pour l'IoT, deux PWA
pour les interfaces et un déploiement sur Cloud IaaS depuis GitHub. Le détail des
technologies est donné dans la section~\ref{sec:techno}.

\subsection{Estimation des charges (CoCoMo+)}
Les charges sont estimées par phase et par rôle, en \textbf{hommes-jours} (h-j).

\begin{tabularx}{\textwidth}{@{}L l c c c c@{}}
\toprule
\tete{Phase} & \tete{Période} & \tete{Front} & \tete{Middleware} & \tete{Back / IoT} & \tete{Total} \\
\midrule
Cahier des charges & jusqu'au 24/10/2026 & 1 & 1 & 1 & 3 \\
Cahier conceptuel & 25/10 -- 07/11/2026 & 3,5 & 4,5 & 4 & 12 \\
Passerelle IoT \& protocoles & 08/11 -- 21/11/2026 & 0 & 2 & 13 & 15 \\
API REST \& middleware & 22/11 -- 12/12/2026 & 0 & 15 & 5 & 20 \\
PWA, intégration \& tests & 13/12/2026 -- 31/01/2027 & 18 & 5 & 2 & 25 \\
\midrule
\textbf{Total} & & \textbf{22,5} & \textbf{27,5} & \textbf{25} & \textbf{75} \\
\bottomrule
\end{tabularx}

\subsection{Planning (diagramme de Gantt)}
\begin{center}
\begin{ganttchart}[
    time slot format=isodate,
    time slot unit=day,
    x unit=0.085cm,
    y unit chart=0.55cm,
    vgrid=false,
    title height=1,
    title label font=\small,
    bar/.style={fill=gray!40,draw=black},
    bar height=0.5,
    bar label font=\small,
    milestone/.style={fill=black},
    milestone label font=\small\itshape
  ]{2026-10-12}{2027-01-31}
  \gantttitlecalendar{month=shortname} \\
  \ganttbar{Cahier des charges}{2026-10-12}{2026-10-24} \\
  \ganttbar{Cahier conceptuel}{2026-10-25}{2026-11-07} \\
  \ganttbar{Passerelle IoT}{2026-11-08}{2026-11-21} \\
  \ganttbar{API REST \& middleware}{2026-11-22}{2026-12-12} \\
  \ganttbar{PWA, intégration \& tests}{2026-12-13}{2027-01-31} \\
  \ganttmilestone{Démonstration finale}{2027-01-31}
\end{ganttchart}
\end{center}

\newpage
% =====================================================================
\section{Diagramme de déploiement}
% =====================================================================

% ---- Pseudo-3D UML node macro ----
\newcommand{\umlnode}[3]{% name, position, content
  \node[node3d] (#1) at #2 {#3};
  \draw[ctblue,thick,fill=ctlight!50]
    (#1.north west) -- ++(0.25,0.25) -- ([xshift=0.25cm,yshift=0.25cm]#1.north east) -- (#1.north east) -- cycle;
  \draw[ctblue,thick,fill=ctlight!50]
    (#1.north east) -- ++(0.25,0.25) -- ([xshift=0.25cm,yshift=0.25cm]#1.south east) -- (#1.south east) -- cycle;
}

\begin{center}
\resizebox{\textwidth}{!}{%
\begin{tikzpicture}[
  font=\footnotesize,
  node3d/.style={draw=ctblue,fill=white,thick,minimum width=#1,inner sep=5pt,align=center},
  node3d/.default=3.6cm,
  env/.style={draw=ctcyan,fill=ctlight,rounded corners=2pt,inner sep=4pt,align=center,minimum width=3.0cm},
  art/.style={draw=ctgrey,fill=white,inner sep=3pt,align=center,font=\scriptsize,minimum width=2.6cm},
  link/.style={thick,ctblue},
  lab/.style={font=\scriptsize\itshape,fill=white,inner sep=1.5pt,text=ctgrey}
]
% Capteurs
\umlnode{esp}{(0,0)}{%
  \guillemotleft device\guillemotright\\\textbf{ESP32 (1..n)}\\[3pt]
  \artbox{\guillemotleft artifact\guillemotright\\firmware.ino}\\[2pt]
  \artbox{Capteurs T° / HR / porte}}

\umlnode{sim}{(0,-4.3)}{%
  \guillemotleft device\guillemotright\\\textbf{Poste de simulation}\\[3pt]
  \artbox{\guillemotleft artifact\guillemotright\\sensor-simulator}}

% Passerelle
\umlnode{gw}{(5.2,-2.1)}{%
  \guillemotleft device\guillemotright\\\textbf{Raspberry Pi}\\\textit{Passerelle IoT}\\[3pt]
  \envbox{\guillemotleft executionEnvironment\guillemotright\\Node-RED}\\[2pt]
  \artbox{\guillemotleft artifact\guillemotright\\flows.json}\\[2pt]
  \artbox{Tampon local}}

% Cloud
\node[draw=ctblue,dashed,thick,rounded corners=4pt,minimum width=5.4cm,minimum height=9.4cm,
      label={[font=\footnotesize\bfseries,text=ctblue]above:\guillemotleft device\guillemotright{} Cloud IaaS — Ubuntu Linux}]
      (cloud) at (11.0,-2.1) {};
\node[env,minimum width=4.4cm] (broker) at (11.0,1.3) {\guillemotleft executionEnvironment\guillemotright\\\textbf{Broker MQTT}\\(Mosquitto, port 8883)};
\node[env,minimum width=4.4cm] (wf) at (11.0,-1.6) {\guillemotleft executionEnvironment\guillemotright\\\textbf{WildFly — Jakarta EE 11}\\[2pt]
  \artbox{\guillemotleft artifact\guillemotright\\coldtrace-api.war}\\[1pt]
  \artbox{PWA statiques (HTTPS 443)}};
\node[env,minimum width=4.4cm] (db) at (11.0,-5.1) {\guillemotleft executionEnvironment\guillemotright\\\textbf{MongoDB}\\[2pt]
  \artbox{mesures \textbullet{} incidents \textbullet{} config}};

% Clients
\umlnode{phone}{(16.6,0.2)}{%
  \guillemotleft device\guillemotright\\\textbf{Smartphone}\\[3pt]
  \artbox{PWA front-office}}
\umlnode{pc}{(16.6,-4.3)}{%
  \guillemotleft device\guillemotright\\\textbf{Poste administrateur}\\[3pt]
  \artbox{PWA back-office}}

% Liens
\draw[link] (esp.east) -- node[lab,above,sloped]{Wi-Fi / MQTT} (gw.west);
\draw[link] (sim.east) -- node[lab,below,sloped]{MQTT} (gw.west);
\draw[link] (gw.east) -- node[lab,above,sloped]{MQTTS (TLS) :8883} (broker.west);
\draw[link] (broker.south) -- node[lab,right]{MQTT (abonnement)} (wf.north);
\draw[link] (wf.south) -- node[lab,right]{driver MongoDB} (db.north);
\draw[link] (wf.east) -- node[lab,above,sloped]{HTTPS / WebSocket} (phone.west);
\draw[link] (wf.east) -- node[lab,below,sloped]{HTTPS :443} (pc.west);
\end{tikzpicture}}%


{\small Figure 1 — Diagramme UML de déploiement de ColdTrace}
\end{center}

Les mesures suivent le chemin capteur $\rightarrow$ passerelle $\rightarrow$ broker MQTT
$\rightarrow$ API Jakarta EE $\rightarrow$ MongoDB, puis sont affichées dans les PWA.
Seul MQTT (avec TLS) est utilisé en version 1 ; la passerelle conserve les mesures en cas
de coupure réseau.

% =====================================================================
\section{Hypothèses}
% =====================================================================

\subsection{Mobilisation des équipes}
\begin{tabularx}{\textwidth}{@{}l l L@{}}
\toprule
\tete{Équipe} & \tete{Période} & \tete{Travail} \\
\midrule
Back-end / IoT & 08/11 -- 21/11/2026 & ESP32, simulateur, Node-RED, broker MQTT \\
Middleware & 22/11 -- 12/12/2026 & API REST, alertes, MongoDB, déploiement cloud \\
Front-end & 13/12/2026 -- 31/01/2027 & PWA front-office et back-office \\
Toute l'équipe & octobre et janvier & Cahiers, tests, documentation, vidéos \\
\bottomrule
\end{tabularx}

\subsection{Hypothèses de déroulement}
\begin{itemize}
  \item des capteurs simulés complètent le matériel disponible ;
  \item une instance Cloud IaaS est accessible pendant toute la durée du projet ;
  \item une courte formation permet au client pilote de participer aux bêta-tests ;
  \item le planning suit les dates du module et les dépendances techniques :
        la passerelle précède l'API, qui précède les PWA.
\end{itemize}

\subsection{Risques}
\begin{tabularx}{\textwidth}{@{}L l L@{}}
\toprule
\tete{Risque} & \tete{Niveau} & \tete{Réponse prévue} \\
\midrule
Matériel indisponible ou en panne & Élevé & Capteurs simulés prêts dès novembre \\
Intégration capteurs $\rightarrow$ cloud difficile & Élevé & Tester le flux complet le plus tôt possible \\
Prise en main de Jakarta EE 11 & Moyen & Travail en binôme, prototype minimal d'abord \\
Accès au Cloud IaaS & Moyen & Provisionner l'instance dès le début \\
Charge des autres modules & Moyen & Marge dans le planning, priorisation des tâches \\
\bottomrule
\end{tabularx}

\end{document}
